\documentclass[a4paper, 12pt]{scrartcl}
\usepackage[ngerman]{babel}
\usepackage[utf8]{inputenc}
\setlength{\parindent}{0em}
\usepackage{eurosym}
\usepackage{amsmath}
\usepackage{amssymb}
\usepackage{dsfont}
\usepackage{polynom}
\usepackage{graphicx}
\usepackage{caption}
\usepackage{tikz}
\usepackage[a4paper,portrait,left=1.0cm,right=1.0cm,top=2cm,bottom=2cm]{geometry}
\usepackage{pgf} % Zur Einbindung der PGF Files
\usepackage{bm}

\usepackage{algorithm}
\usepackage{algpseudocode}  % pseudocode-Umgebung (\begin{algorithmic}, \State, \For, ...)


% Light Mode Customization
\usepackage{xcolor}
\pagecolor{white}  % Hintergrund auf Weiß
\color{black}      % Text auf Schwarz

% Eigene Farben definiert (nun optimiert für hellen Hintergrund)
\definecolor{rot}{RGB}{204,0,0}
\definecolor{blau}{RGB}{0,102,204}
\definecolor{orange}{RGB}{204,85,0}
\definecolor{gruen}{RGB}{0,153,0}
\definecolor{dunkelgruen}{RGB}{0,153,0}
\definecolor{dunkelblau}{RGB}{0,128,128}
\definecolor{backcolor}{RGB}{235,245,255}
\definecolor{frontcolor}{RGB}{0,153,0}
\definecolor{hellblau}{RGB}{180,210,235}
\definecolor{lila}{RGB}{153,0,153}
\definecolor{gelb}{RGB}{204,153,0}
\definecolor{dunkelrot}{RGB}{153,0,51}
\definecolor{dunkelgelb}{RGB}{153,153,0}
\definecolor{boxgray}{RGB}{245,245,245}        % Sehr helles Grau für Box-Hintergründe
\definecolor{boxframe}{RGB}{200,200,200}       % Helles Grau für Box-Rahmen

\usepackage{footnotehyper}
\BeforeBeginEnvironment{tcolorbox}{\savenotes}
\AfterEndEnvironment{tcolorbox}{\spewnotes}

\usepackage{tcolorbox}
\tcbuselibrary{skins, breakable}

\newcounter{satzcnt}[section]
\renewcommand{\thesatzcnt}{\thesection.\arabic{satzcnt}}

\NewDocumentEnvironment{satz}{o}{%
    \begin{tcolorbox}[
        enhanced, frame hidden, breakable, sharp corners,
        colback=white,
        borderline east={1pt}{0pt}{rot},
        left=-3pt, right=10pt, top=6pt, bottom=6pt,
    ]%
    \refstepcounter{satzcnt}\textbf{Satz~\thesatzcnt\IfValueT{#1}{: #1}.}%
}{%
    \end{tcolorbox}%
}

\NewDocumentEnvironment{definition}{o}{%
    \begin{tcolorbox}[
        enhanced, frame hidden, breakable, sharp corners,
        colback=white,
        borderline east={1pt}{0pt}{gelb},
        left=-3pt, right=10pt, top=6pt, bottom=6pt,
    ]%
    \refstepcounter{satzcnt}\textbf{Definition~\thesatzcnt\IfValueT{#1}{: #1}.}\quad%
}{%
    \end{tcolorbox}%
}

\NewDocumentEnvironment{beispiel}{o}{%
    \begin{tcolorbox}[
        enhanced, frame hidden, breakable, sharp corners,
        colback=white,
        borderline east={1pt}{0pt}{lila},
        left=-3pt, right=10pt, top=6pt, bottom=6pt,
    ]%
    \refstepcounter{satzcnt}\textbf{Beispiel~\thesatzcnt\IfValueT{#1}{: #1}.}\quad%
}{%
    \end{tcolorbox}%
}

\NewDocumentEnvironment{bem}{o}{%
    \begin{tcolorbox}[
        enhanced, frame hidden, breakable, sharp corners,
        colback=white,
        %borderline east={1pt}{0pt}{lila},
        left=-3pt, right=10pt, top=6pt, bottom=6pt,
    ]%
    \refstepcounter{satzcnt}\textbf{Korollar \thesatzcnt\IfValueT{#1}{: #1}.}%\quad%
}{%
    \end{tcolorbox}%
}

\NewDocumentEnvironment{beweis}{}{%
    \begin{tcolorbox}[
        enhanced, frame hidden,breakable, sharp corners,
        colback=white,
        borderline east={1pt}{0pt}{gruen},
        left=-3pt, right=10pt, top=6pt, bottom=6pt,
    ]%
    \textbf{Beweis:}\quad%
}{%
    \hfill$\Box$%
    \end{tcolorbox}%
}


% Listings für Python
%\usepackage[german,tworuled]{algorithm2e}
\usepackage{listings}
\lstset{
    commentstyle=\color{gruen},
    keywordstyle=\color{blau},
    language=Python,
    backgroundcolor=\color{boxgray},
    basicstyle=\ttfamily\color{black},
    stringstyle=\color{rot},
    tabsize=2,
    literate={@}{{@}}1,
    title=\lstname
}

% Hyperref Settings
\usepackage[pdfstartview=Fit,plainpages=false,colorlinks=true,linkcolor=gruen,linktocpage=true,hyperfootnotes=false,bookmarksopen=false]{hyperref}

% Eigene Kommandos
\newcommand{\e}[1]{\text{e}^{#1}}
\newcommand{\rot}[2][rot]{\textcolor{#1}{#2}}
\newcommand{\gruen}[2][gruen]{\textcolor{#1}{#2}}
\newcommand{\frontcolor}[2][frontcolor]{\textcolor{#1}{#2}}
\newcommand{\gelb}[2][gelb]{\textcolor{#1}{#2}}
\newcommand{\orange}[2][orange]{\textcolor{#1}{#2}}
\newcommand{\blau}[2][blau]{\textcolor{#1}{#2}}
\newcommand{\hellblau}[2][hellblau]{\textcolor{#1}{#2}}
\newcommand{\lila}[2][lila]{\textcolor{#1}{#2}}
\newcommand{\dunkelrot}[2][dunkelrot]{\textcolor{#1}{#2}}
\newcommand{\dunkelgelb}[2][dunkelgelb]{\textcolor{#1}{#2}}


\newcommand{\norm}[1]{\left\lVert#1\right\lVert}
\newcommand{\dv}{\thinspace \mathsf{d}v}
\newcommand{\dw}{\thinspace \mathsf{d}w}
\newcommand{\dx}{\thinspace \mathsf{dx}}
\newcommand{\dy}{\thinspace \mathsf{d}y}
\newcommand{\dt}{\thinspace \mathsf{dt}}
\newcommand{\dr}{\thinspace \mathsf{d}r}
\newcommand{\ds}{\thinspace \mathsf{ds}}
\newcommand{\du}{\thinspace \mathsf{du}}
\newcommand{\dz}{\thinspace \mathsf{d}z}
\newcommand{\dH}{\thinspace \mathsf{d}H}
\newcommand{\dS}{\thinspace \mathsf{d}S}
\newcommand{\dW}{\thinspace \mathsf{d}W}
\newcommand{\dX}{\thinspace \mathsf{d}X}
\newcommand{\dY}{\thinspace \mathsf{d}Y}
\newcommand{\dZ}{\thinspace \mathsf{d}Z}
\newcommand{\dP}{\thinspace \mathsf{d}P}
\DeclareMathOperator*{\esssup}{ess\,sup}
\DeclareMathOperator*{\argmin}{arg\,min}
\usepackage[backend=biber,style=numeric]{biblatex}
\addbibresource{references.bib}
\definecolor{codebg}{rgb}{0.92,0.92,0.92}
\definecolor{keywordcolor}{rgb}{0.2,0.2,0.7}
\definecolor{stringcolor}{rgb}{0.8,0.1,0.1}
\definecolor{commentcolor}{rgb}{0.4,0.4,0.4}


\begin{document}

\section{Mathematische Formulierung}

\begin{satz}%[Standardform]
    \label{satz:christoffel}
    Sei $L \in C^2$ aus \ref{satz:ELE} gegeben durch
    %
    \begin{align*}
        L(\bm{q},\bm{\dot{q}}) 
        = 
        T(\bm{q},\bm{\dot{q}}) - V(\bm{q})
        =
        \frac{1}{2} \bm{\dot{q}}^T M(\bm{q}) \bm{\dot{q}} - V(\bm{q})
    \end{align*}
    
    mit $M \in C^2(\mathbb{R}^n; \mathbb{R}^{n \times n})$ symmetrisch und
    positiv definit, $V \in C^2(\mathbb{R}^n; \mathbb{R})$. 
    Dann ist die notwendige Bedingung aus \ref{satz:ELE} äquivalent zu
    %
    \begin{align*}
        M(\bm{q}) \ddot{\bm{q}}
        +
        C(\bm{q},\bm{\dot{q}}) \bm{\dot{q}}
        +
        G(\bm{q})
        =
        \bm{u}(t)
    \end{align*}
    mit den \emph{Christoffel-Symbolen}
    %
    \begin{align*}
        \Gamma_{ijk}(\bm{q})
        &:=
        \frac{1}{2}
        \left(
            \frac{\partial M_{ij}}{\partial q_k}(\bm{q})
            +
            \frac{\partial M_{ik}}{\partial q_j}(\bm{q})
            -
            \frac{\partial M_{kj}}{\partial q_i}(\bm{q})
        \right),
        \\
        C_{ij}(\bm{q},\bm{\dot{q}})
        &:=
        \sum_{k=1}^n \Gamma_{ijk}(\bm{q}) \dot{q}_k,
        \\
        G_i(\bm{q})
        &:=
        \frac{\partial V}{\partial q_i}(\bm{q}).
    \end{align*}
\end{satz}

\begin{beweis}
    Es gilt $\frac{\partial L}{\partial \bm{\dot{q}}} = \frac{\partial T}{\partial \bm{\dot{q}}}$
    und wir berechnen die Ableitung von $T$ komponentenweise
    % dT/dq'
    \begin{align*}
        \frac{\partial T}{\partial \dot{q}_i}
        =
        \frac{\partial}{\partial \dot{q}_i}
        \left(
            \frac{1}{2} 
            \sum_{j=1}^{n}
            \sum_{k=1}^{n} \dot{q}_k M_{kj}(\bm{q}) \dot{q}_j
        \right)
        =
        \frac{1}{2} 
        \sum_{j=1}^{n}
        \sum_{k=1}^{n} M_{kj}(\bm{q})
        \cdot 
        \underbrace{
            \frac{\partial}{\partial \dot{q}_i} 
            \left(
                \dot{q}_k  \dot{q}_j
            \right)
        }_{
            \delta_{ik} \dot{q}_j +  \delta_{ij} \dot{q}_k
        }
        =
        \frac{1}{2} \sum_{j=1}^{n}
        M_{ij}(\bm{q}) \dot{q}_j 
        +
        \frac{1}{2}  
        \sum_{k=1}^{n} M_{ki}(\bm{q}) \dot{q}_k
    \end{align*}
    %
    Da $M$ symmetrisch ist $M_{ij} = M_{ji}$ gilt mit dem Indexshift $k=j$:
    %
    \begin{align*}
        \frac{\partial T}{\partial \dot{q}_i}
        =
        \frac{1}{2} \sum_{j=1}^{n}
        M_{ij}(\bm{q}) \dot{q}_j 
        +
        \frac{1}{2}  
        \sum_{j=1}^{n} M_{ji}(\bm{q}) \dot{q}_j
        =
        \sum_{j=1}^{n} M_{ij}(\bm{q}) \dot{q}_j
    \end{align*}
    %
    Die Zeitableitung ist nach der Produktregel:
    %
    \begin{align*}
        \frac{\mathsf{d}}{\mathsf{d}t}
        \frac{\partial L}{\partial \dot{q}_i}
        =
        \frac{\mathsf{d}}{\mathsf{d}t}
        \frac{\partial T}{\partial \dot{q}_i}
        =
        \sum_{j=1}^{n} M_{ij}(\bm{q}) \ddot{q}_j
        +
        \sum_{j=1}^{n}
        \underbrace{
            \sum_{k=1}^{n} 
            \frac{\partial M_{ij}(\bm{q})}{\partial q_k} \cdot \dot{q}_k
        }_{
            \frac{\mathsf{d}M_{ij}(\bm{q}(t))}{\mathsf{d}t}
        }
        \cdot
        \dot{q}_j
    \end{align*}
    %
    Weiter ist
    %
    \begin{align*}
        \frac{\partial L}{\partial q_i}
        =
        \frac{\partial}{\partial q_i}
        \left(
            \frac{1}{2} 
            \sum_{j=1}^{n}
            \sum_{k=1}^{n} \dot{q}_k M_{kj}(\bm{q}) \dot{q}_j
        \right)
        -
        \frac{\partial V(\bm{q})}{\partial q_i}
        =
        \frac{1}{2} 
        \sum_{j=1}^{n}
        \sum_{k=1}^{n} 
        \frac{\partial M_{kj}(\bm{q})}{\partial q_i} 
        \dot{q}_k 
        \dot{q}_j
        -
        \frac{\partial V(\bm{q})}{\partial q_i}.
    \end{align*}

    Eingesetzt in Gleichung (\ref{eq:ELE_control})
    %
    \begin{align*}
        \sum_{j=1}^n M_{ij}(\bm{q}) \ddot{q}_j
        +
        \sum_{j=1}^{n}
        \sum_{k=1}^{n} 
        \left(
            \frac{\partial M_{ij}(\bm{q})}{\partial q_k} 
            -
            \frac{1}{2} 
            \cdot 
            \frac{\partial M_{kj}(\bm{q})}{\partial q_i} 
        \right)
        \cdot \dot{q}_k \cdot \dot{q}_j
        +
        \frac{\partial V(\bm{q})}{\partial q_i}
        =
        u_i.
    \end{align*}

    Wir schreiben den Term um:
    %
    \begin{align*}
        \frac{1}{2}
        \sum_{j=1}^{n}
        \sum_{k=1}^{n} 
        \frac{\partial M_{ij}(\bm{q})}{\partial q_k} 
        \cdot \dot{q}_k \cdot \dot{q}_j
        =
        \frac{1}{2}
        \sum_{j=1}^{n}
        \sum_{k=1}^{n} 
        \frac{\partial M_{ik}(\bm{q})}{\partial q_j} 
        \cdot \dot{q}_j \cdot \dot{q}_k
    \end{align*}
    %
    und erhalten abschließend:
    %
    \begin{align*}
        \sum_{j=1}^n M_{ij}(\bm{q}) \ddot{q}_j
        +
        \sum_{j=1}^{n}
        \underbrace{
            \sum_{k=1}^{n} 
            \underbrace{
                \frac{1}{2} 
                \left(
                    \frac{\partial M_{ij}(\bm{q})}{\partial q_k} 
                    +
                    \frac{\partial M_{ik}(\bm{q})}{\partial q_j}
                    -
                    \frac{\partial M_{kj}(\bm{q})}{\partial q_i} 
                \right)
            }_{
                =: \Gamma_{ijk}(\bm{q})
            }
            \cdot \dot{q}_k 
        }_{
            C_{ij}(\bm{q},\bm{\dot{q}}) 
        }
        \cdot \dot{q}_j
        +
        \underbrace{
            \frac{\partial V(\bm{q})}{\partial q_i}
        }_{
            G_i(\bm{q})
        }
        =
        u_i.
    \end{align*}
\end{beweis}

\section{Two Link Revolute Manipulator}

Wir bezeichnen $\textbf{q} = 
\begin{bmatrix} \gruen{\theta_1} \\ \dunkelrot{\theta_2} \end{bmatrix}$ 
als die generalisierten Koordinaten des Systems. Die Schwerpunkte der beiden Links
bezeichnen wir mit $(x_1, y_1)$ und $(x_2, y_2)$. Die Abstände der Schwerpunkte
von den Gelenken werden mit $a_1$ und $a_2$ bezeichnet.

% Graphik des Roboterarms
%--------------------------------------------------------------------------------------------------------
\begin{center}
\begin{tikzpicture}[scale = 1.75, every node/.style={transform shape}]
    % Koordinatensystem
    %\draw[very thin, gray] (-5,-5) grid (5,5); % Raster
    %\draw[->] (-5,0) -- (5,0); % x-Achse
    %\draw[->] (0,-5) -- (0,5); % y-Achse
    
    % Achsen und Hilfslinien
    \draw[white, dashed, ->] (0.0,-1.0) -- (3.5,-1.0);
    \draw[white, dashed, ->] (0.0,-1.0) -- (0.0,3.0);
    
    % Schwerpunkte
    \draw[dashed] (1.0,-0.5) -- (1.0,-1.0);
    \draw[dashed] (1.0,-0.5) -- (0,-0.5);
    \node at (-0.2, -0.5) {\tiny $y_1$};  % Punktbeschriftung
    \node at (1, -1.2) {\tiny $x_1$}; 
    \draw[dashed] (2.5,1) -- (0.0,1);
    \draw[dashed] (2.5,1) -- (2.5,-1);
    \node at (-0.2, 1.0) {\tiny $y_2$};  % Punktbeschriftung
    \node at (2.5, -1.2) {\tiny $x_2$}; 

    % Erste Zeichnung (Original)
    \fill[hellblau] (-4,-1) circle (0.25);  % Füllt das Objekt mit hellblau
    \draw[white, line width = 0.2mm] (-4,-1) circle (0.25);  % Weißer Rand um das Objekt
    
    \draw[white, line width = 0.2mm] (-4.25,-1.4) rectangle (-3.75,-1);  % Weißer Rand um das Rechteck
    \fill[hellblau] (-4.25,-1.4) rectangle (-3.75,-1);  % Füllt das Rechteck mit hellblau
    \draw[thick, white] (-4.5,-1.4) -- (-3.5,-1.4);
    
    % Grüne Linie angepasst (Endpunkt auf (-2.0, 0.0))
    \draw[thick, gruen] (-4,-1) -- (-2.0,0.0);
    
    % Rote Linie angepasst (von (-2.0, 0.0) nach (-1.0, 2.0))
    \draw[thick, dunkelrot] (-2.0,0.0) -- (-1.0,2.0);
    
    % Glieder
    \fill[blau] (-2.0,0.0) circle (0.1);
    \draw[white, thin] (-2.0,0.0) circle (0.1);  % Weißer Rand um das Objekt
    \fill[blau] (-4.0,-1.0) circle (0.1);
    \draw[white, thin] (-4.0,-1.0) circle (0.1);  % Weißer Rand um das Objekt
    
	% Beschriftungen
	\node at (-3, -0.75) {\tiny $\gruen{l_1}$};  % 
	\node at (-1.5, 1.5) {\tiny $\dunkelrot{l_2}$};  % 
	
	% Schwerpunkte 
	\fill[white] (-3.0,-0.5) circle (0.05);
	\fill[white] (-1.5,1) circle (0.05);
	
	\draw[white, <->] (-1.85,0.15) -- (-1.5,0.85);
	\draw[white, <->] (-3.85,-0.85) -- (-3.1,-0.45);
	
	% Beschriftungen
	\node at (-3.5, -0.4) {\tiny $a_1$};  % 
	\node at (-1.5, 0.5) {\tiny $a_2$};  % 
    
    % Zweite Zeichnung (verschoben um 4 Einheiten entlang der x-Achse)
    \fill[hellblau] (0,-1) circle (0.25);  % Füllt das Objekt mit hellblau
    \draw[white, line width = 0.2mm] (0,-1) circle (0.25);  % Weißer Rand um das Objekt
    
    \draw[white, line width = 0.2mm] (-0.25,-1.4) rectangle (0.25,-1);  % Weißer Rand um das Rechteck
    \fill[hellblau] (-0.25,-1.4) rectangle (0.25,-1);  % Füllt das Rechteck mit hellblau
    \draw[thick, white] (-0.5,-1.4) -- (0.5,-1.4);
    
    % Grüne Linie (verschoben)
    \draw[thick, gruen] (0,-1) -- (2.0,0.0);
    \draw[thick, gruen, dotted] (0,-1) -- (3.5,0.75);
    
    % Rote Linie (verschoben)
    \draw[thick, dunkelrot] (2.0,0.0) -- (3.0,2.0);
    
    % Glieder
    \fill[blau] (2.0,0.0) circle (0.1);
    \draw[white, thin] (2.0,0.0) circle (0.1);  % Weißer Rand um das Objekt
    \fill[blau] (0.0,-1.0) circle (0.1);
    \draw[white, thin] (0.0,-1.0) circle (0.1);  % Weißer Rand um das Objekt
    \fill[blau] (3.0,2.0) circle (0.1);
    \draw[white, thin] (3.0,2.0) circle (0.1);  % Weißer Rand um das Objekt
    
    % Winkel
    \draw [gruen, thick] (1,-1)  arc[start angle=0,end angle=45,radius=0.5cm];
	\node at (0.75, -0.85) {\tiny $\gruen{\theta_1}$};  % Winkelbeschriftung
	
	\draw [dunkelrot, thick] (3,0.5)  arc[start angle=0,end angle=80,radius=0.5cm];
	\node at (2.7, 0.6) {\tiny $\dunkelrot{\theta_2}$};  % Winkelbeschriftung
\end{tikzpicture}
\end{center}
%--------------------------------------------------------------------------------------------------------

Für die Koordinaten der Schwerpunkte gilt:
% Koordinaten der Schwerpunkte
\begin{align*}
    \begin{pmatrix}
        x_1 \\
        y_1 \\
        z_1
    \end{pmatrix}
    =
    \begin{pmatrix}
        a_1 \cdot \cos(\gruen{\theta_1}) \\
        a_1 \cdot \sin(\gruen{\theta_1}) \\
        0
    \end{pmatrix}
    \quad
    \begin{pmatrix}
        x_2 \\
        y_2 \\  
        z_2
    \end{pmatrix}
    =
    \begin{pmatrix}
        l_1 \cdot \cos(\gruen{\theta_1}) 
        + a_2 \cos(\gruen{\theta_1} + \dunkelrot{\theta_2}) \\
        l_1 \cdot \sin(\gruen{\theta_1}) 
        + a_2 \sin(\gruen{\theta_1} + \dunkelrot{\theta_2}) \\
        0
    \end{pmatrix}
\end{align*}

\subsubsection{Formulierung der Massenmatrix}

Für die Massenmatrix gilt hier:
% M(q)
\begin{align*}
    M(\textbf{q}) 
    = m_1 J_{v_{\gruen{1}}}^T J_{v_{\gruen{1}}} 
    + J_{\omega_{\gruen{1}}}^T I_1 J_{\omega_{\gruen{1}}}  
    + m_2 J_{v_{\dunkelrot{2}}}^T J_{v_{\dunkelrot{2}}} 
    + J_{\omega_{\dunkelrot{2}}}^T I_2 J_{\omega_{\dunkelrot{2}}} 
\end{align*}

Hierbei bezeichnet $J$ die Jacobi-Matrix, $I_i$ das Trägheitstensor 
und $m_i$ die Masse des $i$-ten Links.
% Berechnung der Jacobi-Matrizen J_v
\begin{align*}
    J_{v_1} &=
    \frac{\partial \textbf{r}_1}{\partial \textbf{q}}
    =
    \frac{\partial (x_1,y_1,z_1)}{\partial (\gruen{\theta_1}, \dunkelrot{\theta_2})}
    =
    \begin{pmatrix}
        \frac{\partial x_1}{\partial \gruen{\theta_1}} 
        & \frac{\partial x_1}{\partial \dunkelrot{\theta_2}} \\
        \frac{\partial y_1}{\partial \gruen{\theta_1}} 
        & \frac{\partial y_1}{\partial \dunkelrot{\theta_2}} \\
        \frac{\partial z_1}{\partial \gruen{\theta_1}} 
        & \frac{\partial z_1}{\partial \dunkelrot{\theta_2}} 
    \end{pmatrix}
    =
    \begin{pmatrix}
        -a_1 \sin(\gruen{\theta_1}) & 0 \\
        a_1 \cos(\gruen{\theta_1}) & 0 \\
        0 & 0
    \end{pmatrix}
    \\
    J_{v_2} &=
    \frac{\partial \textbf{r}_2}{\partial \textbf{q}}
    =
    \frac{\partial (x_2,y_2,z_2)}{\partial (\gruen{\theta_1}, \dunkelrot{\theta_2})}
    =
    \begin{pmatrix}
        \frac{\partial x_2}{\partial \gruen{\theta_1}} 
        & \frac{\partial x_2}{\partial \dunkelrot{\theta_2}} \\
        \frac{\partial y_2}{\partial \gruen{\theta_1}} 
        & \frac{\partial y_2}{\partial \dunkelrot{\theta_2}} \\
        \frac{\partial z_2}{\partial \gruen{\theta_1}}
        & \frac{\partial z_2}{\partial \dunkelrot{\theta_2}}
    \end{pmatrix} 
    =
    \begin{pmatrix}
        -l_1 \sin(\gruen{\theta_1}) - a_2 \sin(\gruen{\theta_1} + \dunkelrot{\theta_2}) 
        & -a_2 \sin(\gruen{\theta_1} + \dunkelrot{\theta_2}) \\
        l_1 \cos(\gruen{\theta_1}) + a_2 \cos(\gruen{\theta_1} + \dunkelrot{\theta_2}) 
        & a_2 \cos(\gruen{\theta_1} + \dunkelrot{\theta_2}) \\
        0 & 0
    \end{pmatrix}
\end{align*}

Für das Matrizenprodukt $J_{v_1}^T J_{v_1}$ gilt:
% Berechnung der Matrizenprodukte J_v1^T * J_v1
\begin{align*}
    J_{v_1}^T J_{v_1} 
    &=
    \begin{pmatrix}
        -a_1 \sin(\gruen{\theta_1}) & a_1 \cos(\gruen{\theta_1}) & 0 \\
        0 & 0 & 0
    \end{pmatrix}
    \cdot
    \begin{pmatrix}
        -a_1 \sin(\gruen{\theta_1}) & 0 \\
        a_1 \cos(\gruen{\theta_1}) & 0 \\
        0 & 0
    \end{pmatrix}
    \\
    &=
    \begin{pmatrix}
        a_1^2 \sin^2(\gruen{\theta_1}) + a_1^2 \cos^2(\gruen{\theta_1}) & 0 \\
        0 & 0
    \end{pmatrix}
    \\
    &=
    \begin{pmatrix}
        a_1^2 & 0 \\
        0 & 0
    \end{pmatrix}
\end{align*}

Für das Matrizenprodukt $J_{v_2}^T J_{v_2}$ gilt:
% Berechnung der Matrizenprodukte J_v2^T * J_v2
\begin{align*}
    J_{v_2}^T J_{v_2} 
    &=
    \begin{pmatrix}
        -l_1 \sin(\gruen{\theta_1}) - a_2 \sin(\gruen{\theta_1} + \dunkelrot{\theta_2}) 
        & l_1 \cos(\gruen{\theta_1}) + a_2 \cos(\gruen{\theta_1} + \dunkelrot{\theta_2}) 
        & 0 \\
        -a_2 \sin(\gruen{\theta_1} + \dunkelrot{\theta_2}) 
        & a_2 \cos(\gruen{\theta_1} + \dunkelrot{\theta_2})
        & 0
    \end{pmatrix}
    \\
    &\quad \cdot
    \begin{pmatrix}
        -l_1 \sin(\gruen{\theta_1}) - a_2 \sin(\gruen{\theta_1} + \dunkelrot{\theta_2}) 
        & -a_2 \sin(\gruen{\theta_1} + \dunkelrot{\theta_2}) \\
        l_1 \cos(\gruen{\theta_1}) + a_2 \cos(\gruen{\theta_1} + \dunkelrot{\theta_2}) 
        & a_2 \cos(\gruen{\theta_1} + \dunkelrot{\theta_2}) \\
        0 & 0
    \end{pmatrix}
    \\
    &=
    \begin{pmatrix}
        (l_1^2 + a_2^2 + 2 l_1 a_2  \cos(\dunkelrot{\theta_{2}})) 
        & (a_2^2 + l_1 a_2  \cos(\dunkelrot{\theta_{2}})) \\
        (a_2^2 + l_1 a_2  \cos(\dunkelrot{\theta_{2}})) 
        & a_2^2
    \end{pmatrix}
\end{align*}

Für die Winkelgeschwindigkeiten setzen wir   
$\bm{\omega}_1 = \begin{pmatrix}0 & 0 & \gruen{\theta_1'}\end{pmatrix}^T$ und 
$\bm{\omega}_2 = \begin{pmatrix} 0 & 0 &\gruen{\theta_1'} + \dunkelrot{\theta_2'}\end{pmatrix}^T$.
% Berechnung der Jacobi-Matrizen J_\omega
\begin{align*}
    J_{\omega_1} &=
    \frac{\partial \bm{\omega}_1}{\partial \dot{\textbf{q}}}
    =
    \frac{\partial (\omega_1)}{\partial (\gruen{\theta_1}', \dunkelrot{\theta_2}')}
    =
    \begin{pmatrix}
        0 & 0 \\
        0 & 0 \\
        \frac{\partial \omega_1}{\partial \gruen{\theta_1}'} 
        & \frac{\partial \omega_1}{\partial \dunkelrot{\theta_2}'} 
    \end{pmatrix}
    =
    \begin{pmatrix}
        0 & 0 \\
        0 & 0 \\
        \gruen{1} & \dunkelrot{0}
    \end{pmatrix}
    \\
    J_{\omega_2} &=
    \frac{\partial \bm{\omega}_2}{\partial \dot{\textbf{q}}}
    =
    \frac{\partial (\omega_2)}{\partial (\gruen{\theta_1}', \dunkelrot{\theta_2}')}
    =
    \begin{pmatrix}
        0 & 0 \\
        0 & 0 \\
        \frac{\partial \omega_2}{\partial \gruen{\theta_1}'} 
        & \frac{\partial \omega_2}{\partial \dunkelrot{\theta_2}'} 
    \end{pmatrix}
    =
    \begin{pmatrix}
        0 & 0 \\
        0 & 0 \\
        \gruen{1} & \dunkelrot{1}
    \end{pmatrix}
\end{align*}

Für den Trägheitstensor $I_i$ gilt:
% Allgemeiner Trägheitstensor
\begin{align*}
    I = 
    \begin{pmatrix}
        I_{xx} & I_{xy} & I_{xz} \\
        I_{yx} & I_{yy} & I_{yz} \\
        I_{zx} & I_{zy} & I_{zz}
    \end{pmatrix}   
    =
    \begin{pmatrix}
        \int (y^2 + z^2) \mathsf{d}m & -\int xy \mathsf{d}m & -\int xz \mathsf{d}m \\
        -\int yx \mathsf{d}m & \int (x^2 + z^2) \mathsf{d}m & -\int yz \mathsf{d}m \\
        -\int zx \mathsf{d}m & -\int zy \mathsf{d}m & \int (x^2 + y^2) \mathsf{d}m
    \end{pmatrix}
\end{align*}

Für symmetrische Körper (z.B. Zylinder, Kugel, Quader) 
gilt: $I_{xy} = I_{xz} = I_{yz} = 0$ also erhalten wir eine Diagonalmatrix. 
% Symmetrischer Trägheitstensor
\begin{align*}
    I_i 
    =
    \begin{pmatrix}
        \int_ y^2 \mathsf{d}m & 0 & 0 \\
        0 & \int x^2 \mathsf{d}m & 0 \\
        0 & 0 & \int (x^2 + y^2) \mathsf{d}m
    \end{pmatrix}
    =
    \begin{pmatrix}
        \int_V \rho y^2 \mathsf{d}V & 0 & 0 \\
        0 & \int_V \rho x^2 \mathsf{d}V & 0 \\
        0 & 0 & \int_V \rho (x^2 + y^2) \mathsf{d}V
    \end{pmatrix}
\end{align*}


Für die Massenmatrix $M(\textbf{q})$ gilt somit:
% Berechnung der Massenmatrix M(q)
\begin{align*}
    M(\textbf{q}) 
    &= m_1 J_{v_{\gruen{1}}}^T J_{v_{\gruen{1}}} 
    + J_{\omega_{\gruen{1}}}^T I_1 J_{\omega_{\gruen{1}}}  
    + m_2 J_{v_{\dunkelrot{2}}}^T J_{v_{\dunkelrot{2}}} 
    + J_{\omega_{\dunkelrot{2}}}^T I_2 J_{\omega_{\dunkelrot{2}}}    
    \\
    &=
    m_1 
    \begin{pmatrix}
        a_1^2 & 0 \\
        0 & 0
    \end{pmatrix}
    +
    \underbrace{
        \begin{pmatrix}
            0 & 0 & \gruen{1} \\
            0 & 0 & \dunkelrot{0}
        \end{pmatrix}
        \cdot
        \begin{pmatrix}
            I_{1,xx} & 0 & 0 \\
            0 & I_{1,yy} & 0 \\
            0 & 0 & I_{1,zz}
        \end{pmatrix}}_{
        = 
        \begin{pmatrix}
            0 & 0 & I_{1,zz} \\
            0 & 0 & 0
        \end{pmatrix}
    }
    \cdot   
    \begin{pmatrix}
        0 & 0 \\
        0 & 0 \\
        \gruen{1} & \dunkelrot{0}
    \end{pmatrix}
    \\
    &\quad +
    m_2
    \begin{pmatrix}
        (l_1^2 + a_2^2 + 2 l_1 a_2  \cos(\dunkelrot{\theta_{2}})) 
        & (a_2^2 + l_1 a_2  \cos(\dunkelrot{\theta_{2}})) \\
        (a_2^2 + l_1 a_2  \cos(\dunkelrot{\theta_{2}})) 
        & a_2^2
    \end{pmatrix}
    \\
    &\quad +
    \underbrace{
        \begin{pmatrix}
            0 & 0 & \gruen{1} \\
            0 & 0 & \dunkelrot{1}
        \end{pmatrix}
        \cdot
        \begin{pmatrix}
            I_{2,xx} & 0 & 0 \\
            0 & I_{2,yy} & 0 \\
            0 & 0 & I_{2,zz}
        \end{pmatrix}}_{
        = 
        \begin{pmatrix}
            0 & 0 & I_{2,zz} \\
            0 & 0 & I_{2,zz}
        \end{pmatrix}
    }   
    \cdot
    \begin{pmatrix}
        0 & 0 \\
        0 & 0 \\
        \gruen{1} & \dunkelrot{1}
    \end{pmatrix}
    \\
    &=
    \begin{pmatrix}
        m_1 a_1^2 + I_{1,zz} + m_2 \cdot (l_1^2 + a_2^2 + 2 l_1 a_2  \cos(\dunkelrot{\theta_{2}})) + I_{2,zz}
        & m_2 (a_2^2 + l_1 a_2  \cos(\dunkelrot{\theta_{2}})) +I_{2,zz}  \\
        m_2 (a_2^2 + l_1 a_2  \cos(\dunkelrot{\theta_{2}})) + I_{2,zz}
        & m_2 a_2^2 +I_{2,zz} 
    \end{pmatrix}
\end{align*}

Die Trägheitselemente $I_{i,xx}$ und $I_{i,yy}$ sind nicht relevant.
Für unser Model verwenden wir den Standardansatz eines 
dünnen Stabes der Länge $l_i$ und Masse $m_i$. 
Somit gilt $z= y \approx 0$ und wir erhalten: $I_{i,xx} = I_{i,yy} =0$
% Tensor langer dünner Stab
\begin{align*}
    I_{i,zz} \approx \int_{-\frac{l_{i}}{2}}^{\frac{l_{i}}{2}} 
    \underbrace{\rho}_{\frac{m_{i}}{l_{i}}} x^2 \mathsf{d}x
    = \frac{1}{3} \cdot \frac{m_{i}}{l_{i}} 
    \cdot 
    \left[ 
        \left( \frac{l_{i}}{2} \right)^3 -  \left( -\frac{l_{i}}{2} \right)^3 
    \right]
    = \frac{1}{12} m_{i} l_{i}^{2}
    \Rightarrow
    I_i     
    = 
    \begin{pmatrix}
        0 & 0 & 0 \\
        0 & 0 & 0 \\
        0 & 0 & \frac{1}{12} m_{i} l_{i}^{2}
    \end{pmatrix}
\end{align*}


\subsubsection{Aufstellen der Coriolis-Matrix}

Aus der Symmetrie der Massenmatrix $M_{jk}=M_{kj}$ folgt direkt für die 
Christoffel-Symbole (\ref{satz:christoffel}) die Symmetrie $\Gamma_{ijk} = \Gamma_{ikj}$. 
Weiter gilt für den Fall $n=2$ explizit:
% Berechnung der Christoffel Symbole:
\begin{align*}
    \Gamma_{111} 
    &=
    \frac{1}{2} 
    \left(
        \frac{\partial M_{11}}{\partial q_1} 
        +
        \frac{\partial M_{11}}{\partial q_1} 
        -
        \frac{\partial M_{11}}{\partial q_1} 
    \right)
    =
    0
    \\
    \Gamma_{121}
    &= 
    \Gamma_{112}
    =
    \frac{1}{2} 
    \left(
        \frac{\partial M_{12}}{\partial q_1} 
        +
        \frac{\partial M_{11}}{\partial q_2} 
        -
        \frac{\partial M_{12}}{\partial q_1} 
    \right)
    =
    -m_2 l_1 a_2 \sin(\theta_2)
    \\
    \Gamma_{122} 
    &=
    \frac{1}{2} 
    \left(
        \frac{\partial M_{12}}{\partial q_2} 
        +
        \frac{\partial M_{12}}{\partial q_2} 
        -
        \frac{\partial M_{22}}{\partial q_1} 
    \right)
    =
    -m_2 l_1 a_2 \sin(\theta_2)
    \\
    \Gamma_{211} 
    &=
    \frac{1}{2} 
    \left(
        \frac{\partial M_{21}}{\partial q_1} 
        +
        \frac{\partial M_{21}}{\partial q_1} 
        -
        \frac{\partial M_{11}}{\partial q_2} 
    \right)
    =
    m_2 l_1 a_2 \sin(\theta_2)
    \\
    \Gamma_{221} 
    &=
    \Gamma_{212} 
    =
    \frac{1}{2} 
    \left(
        \frac{\partial M_{22}}{\partial q_1} 
        +
        \frac{\partial M_{21}}{\partial q_2} 
        -
        \frac{\partial M_{12}}{\partial q_2} 
    \right)
    =
    0
    \\
    \Gamma_{222} 
    &=
    \frac{1}{2} 
    \left(
        \frac{\partial M_{22}}{\partial q_2} 
        +
        \frac{\partial M_{22}}{\partial q_2} 
        -
        \frac{\partial M_{22}}{\partial q_2} 
    \right)
    =
    0 
\end{align*}

Für die partiellen Ableitungen der kinetischen Energie 
$T(\textbf{q}, \dot{\textbf{q}}) = \frac{1}{2} \dot{\textbf{q}}^T M(\textbf{q}) \dot{\textbf{q}}$ gilt:
% Ableitung dT/d theta
\begin{align*}
    \frac{\partial T}{\partial \gruen{\theta_{1}}}
    &= 
    \frac{1}{2}  \frac{\partial M_{11}}{\partial \gruen{\theta_{1}}} 
    \gruen{\theta_{1}'}^2 
    + \frac{\partial M_{12}}{\partial \gruen{\theta_{1}}} \gruen{\theta_{1}'} 
    \dunkelrot{\theta_{2}'} 
    + \frac{1}{2} \frac{\partial M_{22}}{\partial \gruen{\theta_{1}}} 
    \dunkelrot{\theta_{2}'}(t)^2
    \\
    \frac{\partial T}{\partial \dunkelrot{\theta_{2}}}
    &= 
    \frac{1}{2} \frac{\partial M_{11}}{\partial \dunkelrot{\theta_{2}}} \gruen{\theta_{1}'}^2
    + \frac{\partial M_{12}}{\partial \dunkelrot{\theta_{2}}} \gruen{\theta_{1}'} \dunkelrot{\theta_{2}'} 
    + \frac{1}{2} \frac{\partial M_{22}}{\partial \dunkelrot{\theta_{2}}} \dunkelrot{\theta_{2}'}(t)^2
    \\
    \frac{\partial T}{\partial \gruen{\theta_{1}'}} 
    &= M_{\gruen{1 1}} \gruen{\theta_{1}'} + M_{\gruen{1}\dunkelrot{2}}  \dunkelrot{\theta_{2}'}
    \\
    \frac{\partial T}{\partial \dunkelrot{\theta_{2}'}} 
    &=
    M_{\gruen{1}\dunkelrot{2}} \gruen{\theta_{1}'} + M_{\dunkelrot{22}} \dunkelrot{\theta_{2}'}
\end{align*}

Für die zeitlichen Ableitungen der Größen gilt 
% Ableitung d/dt dT/d theta_1'
\begin{align*}
    \frac{\mathsf{d}}{\dt} \frac{\partial T}{\partial \gruen{\theta_{1}'}} 
    &=
    \underbrace{
        \frac{\mathsf{d}}{\dt} M_{\gruen{1 1}}
    }_{
        \frac{\partial M_{11}}{\partial \theta_{1}} \theta_{1}' 
        + \frac{\partial M_{11}}{\partial \theta_{2}} \theta_{2}'
        } 
    \cdot \gruen{\theta_{1}'} + M_{\gruen{1 1}} \gruen{\theta_{1}''} + 
    \underbrace{
        \frac{\mathsf{d}}{\dt} M_{\gruen{1}\dunkelrot{2}}
        }_{
            \frac{\partial M_{12}}{\partial \theta_{1}} \theta_{1}' + \frac{\partial M_{12}}{\partial \theta_{2}} \theta_{2}'
            } 
        \cdot \dunkelrot{\theta_{2}'} + M_{\gruen{1}\dunkelrot{2}} \dunkelrot{\theta_{2}''} 
    \\
    &= \frac{\partial M_{11}}{\partial \gruen{\theta_{1}}} \gruen{\theta_{1}'}^2
    + \frac{\partial M_{11}}{\partial \dunkelrot{\theta_{2}}} \gruen{\theta_{1}'} \dunkelrot{ \theta_{2}'} 
    +  M_{\gruen{1 1}} \gruen{\theta_{1}''} 
    + \frac{\partial M_{12}}{\partial \gruen{\theta_{1}}} \gruen{\theta_{1}'} \dunkelrot{\theta_{2}'} 
    + \frac{\partial M_{12}}{\partial \dunkelrot{\theta_{2}}} \dunkelrot{\theta_{2}'}^2 
    + M_{\gruen{1}\dunkelrot{2}} \dunkelrot{\theta_{2}''} \\
    &=
    \frac{\partial M_{11}}{\partial \gruen{\theta_{1}}} \gruen{\theta_{1}'}^2
    + \left[ 
        \frac{\partial M_{11}}{\partial \dunkelrot{\theta_{2}}} + \frac{\partial M_{12}}{\partial \gruen{\theta_{1}}} 
        \right] 
    \gruen{\theta_{1}'} \dunkelrot{ \theta_{2}'} 
    + \frac{\partial M_{12}}{\partial \dunkelrot{\theta_{2}}} \dunkelrot{\theta_{2}'}^2 
    + M_{\gruen{1 1}} \gruen{\theta_{1}''} + M_{\gruen{1}\dunkelrot{2}} \dunkelrot{\theta_{2}''}
\end{align*}

% Ableitung d/dt dT/d theta_2'
\begin{align*}
\frac{\mathsf{d}}{\dt} \frac{\partial T}{\partial \dunkelrot{\theta_{2}'}} 
&=
\underbrace{\frac{\mathsf{d}}{\dt} M_{\gruen{1}\dunkelrot{2}}}_{\frac{\partial M_{12}}{\partial \theta_{1}} \theta_{1}' + \frac{\partial M_{12}}{\partial \theta_{2}} \theta_{2}'} \cdot \gruen{\theta_{1}'} + M_{\gruen{1}\dunkelrot{2}} \gruen{\theta_{1}''} + \underbrace{\frac{\mathsf{d}}{\dt} M_{\dunkelrot{22}}}_{\frac{\partial M_{22}}{\partial \theta_{1}} \theta_{1}' + \frac{\partial M_{22}}{\partial \theta_{2}} \theta_{2}'} \cdot \dunkelrot{\theta_{2}'} + M_{\dunkelrot{22}} \dunkelrot{\theta_{2}''} \\
&= 
\frac{\partial M_{12}}{\partial \theta_{1}}  \gruen{\theta_{1}'}^2 
 + \frac{\partial M_{12}}{\partial \theta_{2}} \gruen{\theta_{1}'} \dunkelrot{\theta_{2}'} 
 + M_{\gruen{1}\dunkelrot{2}} \gruen{\theta_{1}''}
 + \frac{\partial M_{22}}{\partial \theta_{1}} \gruen{\theta_{1}'} \dunkelrot{\theta_{2}'} 
 + \frac{\partial M_{22}}{\partial \theta_{2}}  \dunkelrot{\theta_{2}'}^2 
 + M_{\dunkelrot{22}} \dunkelrot{\theta_{2}''}
\end{align*}

Es gilt 
% Ableitung d/dt dT/d theta_1' - dT/d theta_1
\begin{align*}
    \frac{\mathsf{d}}{\dt} \frac{\partial T}{\partial \gruen{\theta_{1}'}}  - \frac{\partial T}{\partial \gruen{\theta_{1}}}
    &=
    \left[ \frac{\partial M_{11}}{\partial \gruen{\theta_{1}}} - \frac{1}{2}  \frac{\partial M_{11}}{\partial \gruen{\theta_{1}}} \right] \gruen{\theta_{1}'}^2
    + \left[ \frac{\partial M_{11}}{\partial \dunkelrot{\theta_{2}}} + \frac{\partial M_{12}}{\partial \gruen{\theta_{1}}} -  \frac{\partial M_{12}}{\partial \gruen{\theta_{1}}} \right] \gruen{\theta_{1}'} \dunkelrot{ \theta_{2}'} 
    + \left[\frac{\partial M_{12}}{\partial \dunkelrot{\theta_{2}}} - \frac{1}{2} \frac{\partial M_{22}}{\partial \gruen{\theta_{1}}} \right] \dunkelrot{\theta_{2}'}^2 \\
    & \quad
    +  M_{\gruen{1 1}} \gruen{\theta_{1}''} 
    + M_{\gruen{1}\dunkelrot{2}} \dunkelrot{\theta_{2}''}
\end{align*}

% Ableitung d/dt dT/d theta_2' - dT/d theta_2
\begin{align*}
\frac{\mathsf{d}}{\dt} \frac{\partial T}{\partial \dunkelrot{\theta_{2}'}} - \frac{\partial T}{\partial \dunkelrot{\theta_{2}}}
&= 
\left[ \frac{\partial M_{12}}{\partial \gruen{\theta_{1}}} - \frac{1}{2} \frac{\partial M_{11}}{\partial \dunkelrot{\theta_{2}}} \right] \gruen{\theta_{1}'}^2 
+ \left[ \frac{\partial M_{12}}{\partial \dunkelrot{\theta_{2}}} + \frac{\partial M_{22}}{\partial \gruen{\theta_{1}}} - \frac{\partial M_{12}}{\partial \dunkelrot{\theta_{2}}} \right] \gruen{\theta_{1}'} \dunkelrot{\theta_{2}'} 
+ \left[ \frac{\partial M_{22}}{\partial \dunkelrot{\theta_{2}}} - \frac{1}{2} \frac{\partial M_{22}}{\partial \dunkelrot{\theta_{2}}} \right]  \dunkelrot{\theta_{2}'}^2 \\
& \quad + M_{\gruen{1}\dunkelrot{2}} \gruen{\theta_{1}''} + D_{\dunkelrot{22}} \dunkelrot{\theta_{2}''} \\
&=
\underbrace{\left[ \frac{\partial M_{12}}{\partial \gruen{\theta_{1}}} \gruen{\theta_{1}'} - \frac{1}{2} \frac{\partial M_{11}}{\partial \dunkelrot{\theta_{2}}} \gruen{\theta_{1}'} \right]}_{C_{\dunkelrot{2}\gruen{1}}} \gruen{\theta_{1}'}
+ \underbrace{\left[ \frac{1}{2} \frac{\partial M_{22}}{\partial \dunkelrot{\theta_{2}}}  \dunkelrot{\theta_{2}'} + \frac{\partial M_{22}}{\partial \gruen{\theta_{1}}} \gruen{\theta_{1}'} \right]}_{C_{\dunkelrot{22}}}  \dunkelrot{\theta_{2}'}
+ M_{\gruen{1}\dunkelrot{2}} \gruen{\theta_{1}''} + M_{\dunkelrot{22}} \dunkelrot{\theta_{2}''}
\end{align*}

Für die Einträge der Matrix $C$ gilt
% Berechnung der Einträge der Coriolis-Matrix C
\begin{align*}
C_{11} 
&=
\frac{1}{2}  \underbrace{\frac{\partial M_{11}}{\partial \gruen{\theta_{1}}} }_{= 0} \gruen{\theta_{1}'} + \underbrace{\frac{\partial M_{11}}{\partial \dunkelrot{\theta_{2}}}}_{ =  -2m_2 l_1 a_2 \sin(\dunkelrot{\theta_2})} \dunkelrot{ \theta_{2}'}
=
-2m_2 l_1 a_2 \sin(\dunkelrot{\theta_2}) \dunkelrot{ \theta_{2}'} 
\\
C_{12}
&=
\underbrace{\frac{\partial M_{12}}{\partial \dunkelrot{\theta_{2}}}}_{= - m_2 l_1 a_2 \sin(\dunkelrot{\theta_2})} \dunkelrot{\theta_{2}'} - \frac{1}{2} \underbrace{\frac{\partial M_{22}}{\partial \gruen{\theta_{1}}}}_{ = 0} \dunkelrot{\theta_{2}'}
=
- m_2 l_1 a_2 \sin(\dunkelrot{\theta_2}) \dunkelrot{\theta_{2}'}
\\
C_{21} 
&= \underbrace{\frac{\partial M_{12}}{\partial \gruen{\theta_{1}}}}_{= 0} \gruen{\theta_{1}'} - \frac{1}{2} \underbrace{\frac{\partial M_{11}}{\partial \dunkelrot{\theta_{2}}}}_{= -2m_2 l_1 a_2 \sin(\dunkelrot{\theta_2}) } \gruen{\theta_{1}'}
=
m_2 l_1 a_2 \sin(\dunkelrot{\theta_2})\gruen{\theta_{1}'}
\\
C_{22} 
&=
\frac{1}{2} \underbrace{\frac{\partial M_{22}}{\partial \dunkelrot{\theta_{2}}}}_{ = 0} \dunkelrot{\theta_{2}'} + \underbrace{ \frac{\partial M_{22}}{\partial \gruen{\theta_{1}}}}_{ = 0} \gruen{\theta_{1}'}
= 0
\end{align*}

Es gilt somit 
% Resultat Coriolis-Matrix C(q, qdot) * qdot
\begin{align*}
C(\textbf{q}, \dot{\textbf{q}}) \cdot \dot{\textbf{q}}
=
\begin{pmatrix}
    -2m_2 l_1 a_2 \sin(\dunkelrot{\theta_2}) \dunkelrot{ \theta_{2}'} 
    &
    - m_2 l_1 a_2 \sin(\dunkelrot{\theta_2}) \dunkelrot{\theta_{2}'} 
    \\
    m_2 l_1 a_2 \sin(\dunkelrot{\theta_2})\gruen{\theta_{1}'}
    &
    0
\end{pmatrix}
\begin{pmatrix}
    \gruen{\theta_{1}'} \\
    \dunkelrot{\theta_{2}'} 
\end{pmatrix}
\end{align*}

\subsubsection{Gravitationsterme}

Für die potenzielle Energie $V$ gilt:
% V = m*g*h als Potenzielle Energie
\begin{align*}
    V 
    &= m_1 \cdot g \cdot y_1 + m_2 \cdot g \cdot y_2 
    = m_1  g  a_1 \sin(\gruen{\theta_1}) 
    + m_2  g \left( l_1 \sin(\gruen{\theta_1}) 
    + a_2 \sin(\gruen{\theta_1} + \dunkelrot{\theta_2}) \right)
\end{align*}

Für die partiellen Ableitungen erhalten wir
% Berechnung der Gravitationsterme G(q) = dV/dq
\begin{align*}
    G(\textbf{q})
    =
    \frac{\partial V}{\partial \textbf{q}}
    =
    \begin{pmatrix}
        \frac{\partial V}{\partial \gruen{\theta_1}} \\
        \frac{\partial V}{\partial \dunkelrot{\theta_2}}
    \end{pmatrix}
    =
    \begin{pmatrix}
        m_1  g  a_1 \cos(\gruen{\theta_1}) + m_2  g \left( l_1 \cos(\gruen{\theta_1}) 
        + a_2 \cos(\gruen{\theta_1} + \dunkelrot{\theta_2}) \right) \\
        m_2  g \left( l_1 \cos(\gruen{\theta_1}) + a_2 \cos(\gruen{\theta_1} 
        + \dunkelrot{\theta_2}) \right)
    \end{pmatrix}
\end{align*}

% Hier Berechnung für zusätzliches Gewicht an link2
\newpage

\subsection{Einfaches Greifobjekt}
Wir modellieren ein Variables Greifobjekt, welches am Ende von Link 2 gehalten wird. 
Wir betrachten nachfolgend das Greifobjekt als Punktmasse mit Masse $m_G$. 

% Greifobjekt
%--------------------------------------------------------------------------------------------------------
\begin{center}
    \begin{tikzpicture}[scale = 1.75, every node/.style={transform shape}]
        % Koordinatensystem
        %\draw[very thin, gray] (-5,-5) grid (5,5); % Raster
        %\draw[->] (-5,0) -- (5,0); % x-Achse
        %\draw[->] (0,-5) -- (0,5); % y-Achse
        
        % Erste Zeichnung (Original)
        \fill[hellblau] (-4,-1) circle (0.25);  % Füllt das Objekt mit hellblau
        %\draw[white, line width = 0.2mm] (-4,-1) circle (0.25);  % Weißer Rand um das Objekt
        
        %\draw[white, line width = 0.2mm] (-4.25,-1.4) rectangle (-3.75,-1);  % Weißer Rand um das Rechteck
        \fill[hellblau] (-4.25,-1.4) rectangle (-3.75,-1);  % Füllt das Rechteck mit hellblau
        %\draw[thick, white] (-4.5,-1.4) -- (-3.5,-1.4);
        
        % Grüne Linie angepasst (Endpunkt auf (-2.0, 0.0))
        \draw[thick, gruen] (-4,-1) -- (-3.0,1.0);
        
        % Rote Linie angepasst (von (-2.0, 0.0) nach (-1.0, 2.0))
        \draw[thick, dunkelrot] (-3.0,1.0) -- (-1.0,1.0);
        
        % Glieder
        \fill[blau] (-3.0,1.0) circle (0.1);
        \draw[white, thin] (-3.0,1.0) circle (0.1);  % Weißer Rand um das Objekt
        \fill[blau] (-4.0,-1.0) circle (0.1);
        \draw[white, thin] (-4.0,-1.0) circle (0.1);  % Weißer Rand um das Objekt
        
        % Beschriftungen
        \node at (-3, -0.75) {\tiny $\gruen{l_1}$};  % 
        \node at (-1.5, 1.5) {\tiny $\dunkelrot{l_2}$};  % 
        
        % Schwerpunkte 
        %\fill[white] (-3.0,-0.5) circle (0.05);
        %\fill[white] (-1.5,1) circle (0.05);
        
        %\draw[white, <->] (-1.85,0.15) -- (-1.5,0.85);
        %\draw[white, <->] (-3.85,-0.85) -- (-3.1,-0.45);
        
        % Beschriftungen
        \node at (-3.5, -0.4) {\tiny $a_1$};  % 
        \node at (-1.5, 0.5) {\tiny $a_2$};  % 

        \draw[dunkelrot, thick] (-1.0,0.9) -- (-1.0, 1.1);
        \draw[dunkelrot, thick] (-1.0,0.9) -- ++(0.1,0.0);
        \draw[dunkelrot, thick] (-1.0,1.1) -- ++(0.1,0.0);
        \node at (-0.6,0.7) {\tiny $m_G$}; % Beschriftung

        % Erste Zeichnung (Original) verschoben um +4 in x
        \fill[hellblau] (0,-1) circle (0.25);  % Füllt das Objekt mit hellblau
        \draw[white, line width = 0.2mm] (0,-1) circle (0.25);  % Weißer Rand um das Objekt
        
        \draw[white, line width = 0.2mm] (-0.25,-1.4) rectangle (0.25,-1);  % Weißer Rand um das Rechteck
        \fill[hellblau] (-0.25,-1.4) rectangle (0.25,-1);  % Füllt das Rechteck mit hellblau
        \draw[thick, white] (-0.5,-1.4) -- (0.5,-1.4);
        
        % Grüne Linie angepasst (Endpunkt auf (1.0, 1.0))
        \draw[thick, gruen] (0,-1) -- (1.0,1.0);
        
        % Rote Linie angepasst (von (1.0, 1.0) nach (3.0, 1.0))
        \draw[thick, dunkelrot] (1.0,1.0) -- (3.0,1.0);

        % Orange Linie zum Greifobjekt
        \draw[thick, orange] (3.0,1.0) -- (3.3,1.4);
        \draw[thick, orange] (3.364,1.352) -- (3.236,1.448);
        \draw[thick, orange] (3.364,1.352) -- ++(0.072,0.096);
        \draw[thick, orange] (3.236,1.448) -- ++(0.072,0.096);

        % Winkel
        \draw [orange, thick] (3.35,1.0)  arc[start angle=0,end angle=45,radius=0.35cm];
	    \node at (3.25, 0.9) {\tiny $\orange{\theta_3}$};  % Winkelbeschriftung
        \draw[thick, dunkelrot, dotted] (3,1) -- (4.0,1.0);
        \node at (3.0, 1.3) {\tiny $\orange{l_3}$};

        \fill[blau] (3.0,1.0) circle (0.1);
        \draw[white, thin] (3.0,1.0) circle (0.1);  % Weißer Rand um das Objekt

        % Glieder
        \fill[blau] (1.0,1.0) circle (0.1);
        \draw[white, thin] (1.0,1.0) circle (0.1);  % Weißer Rand um das Objekt
        \fill[blau] (0.0,-1.0) circle (0.1);
        \draw[white, thin] (0.0,-1.0) circle (0.1);  % Weißer Rand um das Objekt
    \end{tikzpicture}
\end{center}
%--------------------------------------------------------------------------------------------------------

Für die Masse von Link 2 inklusive Greifobjekt gilt:
% m_i' = m_i + m_G
\begin{align*}
    m_{i}' = m_{i} + m_{G}
\end{align*}

Für den Schwerpunkt von Link 2 inklusive Greifobjekt gilt:
% a_i' = (m_i * a_i + m_G * d) / (m_i + m_G)
\begin{align*}
    \textbf{a}_{2}'
    &= \frac{m_{i} \cdot \textbf{a}_{2} + m_{G} \cdot \textbf{d}}{m_{i} + m_{G}}
    \quad
    \textbf{d}(\orange{\theta_3}) =
    \begin{pmatrix}
        l_2 + r \cdot \cos(\orange{\theta_3}) \\
        r \cdot \sin(\orange{\theta_3}) \\
        0
    \end{pmatrix}
\end{align*}

Für den Trägheitstensor von Link 2 inklusive Greifobjekt gilt:
% I_i' = I_i + m_i * S(a_i - a_i') + m_G * S(d - a_i')
\begin{align*}
    I_{i}' 
    &= 
    I_{i} + m_{i} \cdot S(\textbf{a}_{i} - \textbf{a}_{i}') 
    + m_{G} \cdot S(\textbf{d} - \textbf{a}_{i}')
    \quad
    S(\textbf{v}) =
    \begin{pmatrix}
        v_{y}^2 + v_{z}^2 & -v_{x} v_{y} & -v_{x} v_{z} \\
        -v_{x} v_{y} & v_{x}^2 + v_{z}^2 & -v_{y} v_{z} \\
        -v_{x} v_{z} & -v_{y} v_{z} & v_{x}^2 + v_{y}^2
    \end{pmatrix}
\end{align*}

Wobei $S(\textbf{v})$ die Steiner-Matrix ist. 
Wir erhalten die Matrix direkt aus der Definition des Trägheitstensors (\ref{eq:traegheitstensor}) einer Punktmasse
mit $\rho(\textbf{r}) = m \cdot \delta(\text{r}-\textbf{d})$.
Für die Verschiebungsvektoren gilt:
% Berechnung a_2'
\begin{align*}
    \textbf{d}
    =
    \begin{pmatrix}
        l_2 \\
        0 \\
        0 
    \end{pmatrix}
    \textbf{a}_{2} =
    \begin{pmatrix}
        \frac{l_2}{2} \\
        0 \\
        0 
    \end{pmatrix}
    \Rightarrow
    \textbf{a}_{2}'
    =
    \frac{l_2}{m_{2} + m_{G}}
    \begin{pmatrix}
        \frac{m_{2}}{2} + m_{G} \\
        0 \\
        0
    \end{pmatrix}
\end{align*}

Für die Steiner-Matrizen gilt entsprechend
% Berechnung der Steiner-Matrizen
\begin{align*}
    S(\textbf{a}_{2} - \textbf{a}_{2}')
    =
    \begin{pmatrix}
        0 & 0 & 0 \\
        0 & \left[\frac{l_2}{2} - a_{2}'\right]^2 & 0 \\
        0 & 0 & \left[\frac{l_2}{2} - a_{2}'\right]^2
    \end{pmatrix}
    \quad
    S(\textbf{d} - \textbf{a}_{2}')
    =
    \begin{pmatrix}
        0 & 0 & 0 \\
        0 & \left[l_2^2 - a_{2}'\right]^2 & 0 \\
        0 & 0 & \left[l_2^2 - a_{2}'\right]^2
    \end{pmatrix}
\end{align*}

Somit erhalten wir für den Trägheitstensor von Link 2 inklusive Greifobjekt:
% Berechnung I_2'
\begin{align*}
    I_{2}'
    &=
    \begin{pmatrix}
        I_{2,xx} & 0 & 0 \\
        0 & I_{2,yy} 
        + m_2 \cdot \left[\frac{l_2^2}{2} - a_{2}'\right]^2 
        + m_G \cdot \left[\frac{l_2^2}{2} - a_{2}'\right]^2 
        & 0 \\
        0 & 0 & I_{2,zz} 
        + m_2 \cdot \left[\frac{l_2^2}{2} - a_{2}'\right]^2 
        + m_G \cdot \left[\frac{l_2^2}{2} - a_{2}'\right]^2
    \end{pmatrix}
\end{align*}

Für die Koordinaten des Schwerpunktes von Link 3 gilt:
% Berechnung der Koordinaten des Schwerpunktes von Link 3
\begin{align*}
    \textbf{r}_3
    &=
    \begin{pmatrix}
        l_1 \cos(\gruen{\theta_1}) 
        + l_2 \cos(\gruen{\theta_1} + \dunkelrot{\theta_2}) 
        + l_3 \cos(\gruen{\theta_1} + \dunkelrot{\theta_2} + \orange{\theta_3}) \\
        l_1 \sin(\gruen{\theta_1}) 
        + l_2 \sin(\gruen{\theta_1} + \dunkelrot{\theta_2}) 
        + l_3 \sin(\gruen{\theta_1} + \dunkelrot{\theta_2} + \orange{\theta_3}) \\
        0
    \end{pmatrix}
\end{align*}
Berechnung der Jacobi-Matrizen für Link 3 bezüglich der Geschwindigkeiten:
% Jacobi-Matrizen für Link 3
\begin{align*}
    J_{v_3} &=
    \frac{\partial \textbf{r}_3}{\partial \textbf{q}}
\end{align*}

Für die Winkelgeschwindigkeit des Links 3 setzen wir 
$\bm{\omega_3}= 
\begin{pmatrix} 
    0 & 0 & \gruen{\theta_1'} + \dunkelrot{\theta_2'} + \orange{\theta_3'} 
\end{pmatrix}$

und erhalten für die Jacobi-Matrix bezüglich der Winkelgeschwindigkeit:
% Jacobi-Matrizen für Link 3
\begin{align*}
    J_{\omega_3} &=
    \frac{\partial \bm{\omega_3}}{\partial \dot{\textbf{q}}} 
    =
    \begin{pmatrix}
        0 & 0 & 0 \\
        0 & 0 & 0 \\
        1 & 1 & 1
    \end{pmatrix}
\end{align*}


\end{document}



%------------------------------------------------------------------------------------------------------------
